\DocumentMetadata{lang=en-US, pdfstandard=ua-2, testphase={phase-III,table}}
\documentclass{./template/articlex}

\hypersetup{pdftitle={Why Economics? What UL Lafayette Graduates Actually Earn},
            pdflang={en-US}}

% --- handout formatting -------------------------------------------------
\usepackage{geometry}
\geometry{landscape, left=0.7in, right=0.7in, top=0.7in, bottom=0.7in}
\usepackage{xcolor}
\usepackage{multicol}
\usepackage{tcolorbox}
\usepackage[scaled=0.92]{helvet}

\definecolor{ULRED}{HTML}{CE0E2D}
\definecolor{INK}{HTML}{222222}
\definecolor{SOFT}{HTML}{555555}

% section heading
\newcommand{\hd}[1]{%
  \vspace{5pt}\noindent{\sffamily\bfseries\color{ULRED}\large #1}\par\vspace{2pt}}

% lead-in phrase at the start of a paragraph
\newcommand{\lead}[1]{\noindent{\sffamily\bfseries\color{INK}#1}}

% one recent-graduate placement
\newcommand{\gradentry}[3]{%
  \noindent{\sffamily\bfseries\color{INK}#1}\\
  {\color{SOFT}#2\quad\textit{Class of #3}}\par\vspace{5pt}}

\setlength{\columnsep}{28pt}
\setlength{\parskip}{3pt}
\setlength{\parindent}{0pt}
% ------------------------------------------------------------------------

\begin{document}
\begin{singlespace}

{\sffamily\bfseries\color{ULRED}\fontsize{26}{30}\selectfont
Why Economics? What UL Lafayette Graduates Actually Earn}

\vspace{2pt}
{\color{ULRED}\rule{\linewidth}{1.5pt}}
\vspace{4pt}

{\sffamily\color{SOFT}\large
If you are picking a major, it helps to know where each one tends to lead. These
figures show what UL Lafayette graduates actually earn five and ten years after
they finish their degree --- and where Economics majors land among them.}

\vspace{6pt}

\begin{tcolorbox}[colback=ULRED!5, colframe=ULRED, boxrule=1pt, arc=2pt,
                  left=10pt, right=10pt, top=5pt, bottom=5pt]
{\sffamily\bfseries\color{ULRED}The bottom line.}\quad
{\color{INK}At five years, the business majors look much alike. By ten years
out, UL Lafayette Economics graduates earn more than graduates of any other
business major here, and about 50\% more than the typical UL graduate across
every field on campus --- at the middle of the pack and at the upper end alike.}
\end{tcolorbox}

\vspace{6pt}

\small
\begin{multicols}{2}

\hd{Where these numbers come from}

The U.S.\ Census Bureau runs a program called \textbf{Post-Secondary Employment
Outcomes}, or PSEO. It takes UL Lafayette's own transcript records and matches
them to the earnings records that employers file for their workers. So these
are the real paychecks of real UL graduates.

This is worth pausing on. \textbf{These are not survey answers.} Nobody mailed
a questionnaire and reported back on whoever felt like replying, which is how a
lot of ``average salary by major'' lists get made. Those lists tend to hear
from the graduates doing well and miss the ones who aren't. Here, the earnings
come straight from payroll records, so there is nobody missing for that reason.

\hd{How to read the figures}

\lead{Median earnings} is the graduate right in the middle. Half of that
major's graduates earned more, half earned less. It is the best single answer
to ``what does a normal graduate make?''

\lead{75th percentile earnings} is the graduate three-quarters of the way up.
One in four graduates from that major earned at least this much. It shows what
things look like for the people doing well in a major, not just the middle of
the pack.

\lead{5 years and 10 years} are counted from the first full year after
graduation. \textbf{Pay attention to the ten-year numbers.} Starting salaries
across majors look a lot more similar than mid-career salaries do, and the
differences that matter for your life show up later.

\lead{January 2026 dollars.} Someone who graduated years ago and someone who
graduated recently are compared fairly, without older salaries looking small
just because prices were lower then.

\hd{Who is counted}

Both figures cover people who earned a \textbf{bachelor's degree at UL
Lafayette between 2001 and 2021}. Census sorts graduates into three-year
classes --- 2001--2003, 2004--2006, on through 2019--2021 --- and these figures
pool all of them together.

To be included, a graduate had to be working steadily that year. Census leaves
out anyone who earned less than a full-time minimum-wage job would pay, and
anyone who went two or more quarters with no earnings. So these are the
earnings of graduates who were actually working.

One wrinkle: the ten-year numbers can only include students who graduated at
least ten years ago, so the ten-year bars describe an earlier set of graduates
than the five-year bars do. Everything is still in January 2026 dollars, so
inflation is not why the ten-year numbers are larger.

\hd{Three things these numbers do not tell you}

\lead{Graduate school is included.} If someone earned a bachelor's in Economics
here and later went on to an MBA, a law degree, or a master's, they still show
up in the Economics numbers. Census does not take them out. Some majors send
more of their graduates on to further study than others, and Economics sends a
lot. So part of what you see is where a major can take you next.

\lead{Double majors are counted twice.} Someone who majored in both Finance and
Economics appears in both bars.

\lead{This is not a promise.} These figures show what graduates of each major
earned, not what the major did for them. Part of the difference is the work
graduates go into, and part is which students pick which major to begin with.

\end{multicols}
\normalsize



\clearpage
\begin{figure}[H]
  \centering
  {\sffamily\bfseries\color{ULRED}\large Figure 1.}
  {\sffamily\large Economics compared with UL Lafayette's other business majors}
  \par\vspace{6pt}
  \includegraphics[width=\linewidth, alt={Bar chart comparing median and 75th-percentile earnings for seven UL Lafayette undergraduate business majors at five and ten years after graduation. At ten years Economics leads both measures, at $105,369 median and $141,919 at the 75th percentile; Accounting is next at $85,829 median.}]{figures/econ_vs_business_majors_5yrs_10yrs_2026-09-28.pdf}
  \label{fig:econ_vs_business}
\end{figure}
\vspace{0pt}
\noindent\begin{minipage}{\linewidth}
\footnotesize{\textbf{Notes:} Each bar is one of UL Lafayette's undergraduate
business majors. Economics is the red bar. The majors are listed in the same
order in all four boxes, sorted by what the middle graduate earns ten years
out, so you can read straight down a column and compare. Figures cover
bachelor's graduates from the classes of 2001 through 2021.
Census reports earnings using national program codes that divide majors more
finely than UL does. Related codes have been grouped into the majors shown here:
Management Information Systems, for example, is counted within Management, and
Taxation within Accounting. Where a major covers more than one code, the figure
shows an average that gives more weight to the codes with more graduates.
Economics is a single code, so its bars are exactly the numbers Census published.
Earnings are in January 2026 dollars, adjusted from Census's published figures
using the Consumer Price Index (BLS series CUUR0000SA0). \textbf{Source:} U.S.
Census Bureau, Longitudinal Employer--Household Dynamics Program,
Post-Secondary Employment Outcomes (PSEO), Louisiana file,
\url{https://lehd.ces.census.gov/data/pseo_experimental.html};
accessed September 28, 2026.}
\end{minipage}


\clearpage
\begin{figure}[H]
  \centering
  {\sffamily\bfseries\color{ULRED}\large Figure 2.}
  {\sffamily\large Economics compared with all UL Lafayette graduates}
  \par\vspace{0pt}
  \includegraphics[width=\linewidth, alt={Bar chart comparing Economics graduates with all UL Lafayette bachelor's graduates at five and ten years after graduation. At ten years Economics earns $105,369 at the median against $70,951 for all majors, and $141,919 at the 75th percentile against $106,301.}]{figures/econ_vs_all_majors_5yrs_10yrs_2026-09-28.pdf}
  \label{fig:econ_vs_all}
\end{figure}
\vspace{-6pt}
\noindent\begin{minipage}{\linewidth}
\footnotesize{\textbf{Notes:} This figure looks past the business school. The
red bar is Economics. The gray bar is every student who earned a bachelor's
degree at UL Lafayette, in any field: nursing, engineering, education, the
sciences, the arts, and business all together. Figures cover bachelor's
graduates from the classes of 2001 through 2021.
That gray bar is not something we calculated. Census publishes a single
university-wide figure for UL Lafayette covering all fields of study, and that
is what is shown, so every program is represented in proportion to its size.
Earnings are in January 2026 dollars, adjusted from Census's published figures
using the Consumer Price Index (BLS series CUUR0000SA0). \textbf{Source:} U.S.
Census Bureau, Longitudinal Employer--Household Dynamics Program,
Post-Secondary Employment Outcomes (PSEO), Louisiana file,
\url{https://lehd.ces.census.gov/data/pseo_experimental.html};
accessed September 28, 2026.}
\end{minipage}


\clearpage
{\sffamily\bfseries\color{ULRED}\fontsize{20}{24}\selectfont
Where Economics Majors Actually Work}

\vspace{2pt}
{\color{ULRED}\rule{\linewidth}{1.5pt}}
\vspace{6pt}

{\sffamily\color{SOFT}\large
An Accounting major applies for jobs called ``Accountant.'' Economics does not
work that way --- and that is the single biggest reason students underestimate
it. Almost none of the jobs below carry the word \emph{economist} in the title.
All of them hire economics majors.}

\vspace{8pt}

\small
\begin{multicols}{2}

\hd{Business and finance}

\lead{Financial Analyst.} Build the forecasts a company uses to decide where its
money goes --- which projects get funded, what a business is worth, what happens
to the bottom line if costs move.

\lead{Credit Analyst / Underwriter.} Decide whether a loan or a policy is a good
risk. You are pricing uncertainty, which is most of what an economics degree
teaches you to do.

\lead{Pricing or Revenue Analyst.} Work out what the market will bear. Hotels,
airlines, energy firms, and hospitals all employ people to do nothing else.

\lead{Management Consultant.} Get handed a business problem you have never seen
before and a pile of data, and asked what is going on and what to do about it.

\hd{Data and research}

\lead{Data Analyst / Business Analyst.} Turn what a company already records into
decisions it can act on. Economics is the major that teaches you to ask whether
a difference in the data is real before you act on it.

\lead{Market Research Analyst.} Measure demand, size a market, test what
customers actually respond to as opposed to what they say they want.

\lead{Research Assistant.} Federal Reserve banks, the Census Bureau, the Bureau
of Labor Statistics, and university research centers hire economics graduates
directly out of a bachelor's degree. It is also the usual route into a PhD.

\hd{Government and policy}

\lead{Policy Analyst.} State agencies, legislative staff, and research
organizations need people who can say what a proposed rule or tax change would
actually do, and at what cost.

\lead{Economic Consultant.} Litigation support, antitrust, and regulatory work:
estimating damages, defining markets, and explaining the estimate to people who
are not economists.

\hd{Around Acadiana}

Energy and commodity analysis, healthcare analytics, insurance and risk,
banking, and port and logistics operations all employ analysts locally, and all
of them are hiring for the same underlying skill.

\hd{What employers are actually buying}

They are not buying a job title. They are buying the ability to work out what is
driving a number, to say what would happen if something changed, and to tell the
difference between a real effect and noise --- then explain it to someone who
does not want to see the math. That skill does not have a single job title,
which is exactly why it travels so far.

Economics is also a well-worn path into law school and into graduate programs in
business, public policy, and economics itself.

\end{multicols}
\normalsize


\begin{comment}
\clearpage
{\sffamily\bfseries\color{ULRED}\fontsize{20}{24}\selectfont
Where Recent UL Economics Graduates Are Working}

\vspace{2pt}
{\color{ULRED}\rule{\linewidth}{1.5pt}}
\vspace{6pt}

{\sffamily\color{SOFT}\large
A sample of placements from recent graduating classes.}

\vspace{10pt}

\small
\begin{multicols}{3}
\input{recent-grads}
\end{multicols}
\normalsize


\end{comment}


\end{singlespace}
\end{document}


